\section{Introduction}
\subsection{Motivation}
This project was initially motivated by my lived experience as a \gls{neurodivergent} developer. In professional settings, \gls{ASD} is frequently 
romanticised as a ``superpower'' due to the high capacity for deep logical analysis
and rapid skill acquisition. However, this perceived strength came with a significant \gls{cognitive_tax}. Adapting to new environments, fluctuating workloads and 
maintaining a sustainable work-life balance is often extremely difficult as current tools like \glspl{IDE} lack out of the box features to mitigate 
these executive functioning barriers. This research seeks to transition the \gls{IDE} from a passive tool into an active cognitive partner through AI-driven support.


\subsection{Problem Statement}
While modern \glspl{IDE} are feature-rich, their design reflects a \gls{neurotypical-bias} that assumes a high baseline for filtering sensory input and managing 
multi-step planning. For developers with \gls{ADHD}, \gls{ASD}, or \gls{Dyslexia}, these tools often exacerbate Extraneous Cognitive Load, 
the mental effort spent on processing the interface rather 
than the problem logic~\cite{sweller1988cognitive}.

Current \glspl{IDE} prioritise feature density, often resulting in visual noise from nested menus and persistent notifications. Research indicates that developers with \gls{ADHD} 
frequently experience \gls{task_paralysis}, difficulty initiating tasks during a project~\cite{bramham2009executive}, while those with \gls{ASD} may struggle with the \gls{double_empathy} problem, leading to them misinterpreting 
ambiguous or implied requirements from team members~\cite{Milton01102012}. Existing AI tools focus almost exclusively on code generation rather than process support. Consequently, there is a critical need 
for an \gls{IDE} intervention that acts as a digital ramp, filtering noise and providing the scaffolding necessary to bypass the executive freeze response.

\subsection{Understanding \Gls{neurodiversity} in Software Development}
Before going into the specific technical barriers of modern \glspl{IDE}, it is necessary to define \gls{neurodiversity} and the cognitive
 frameworks that underpin this research. Originally coined by Judy Singer in the late 1990s, \gls{neurodiversity} views neurological differences, such 
 as \gls{ASD}, \gls{ADHD}, and \gls{Dyslexia}, as natural variations in the human genome rather than purely medical deficits~\cite{singer1999why}.

In the context of software engineering, these variations often manifest as unique cognitive processing styles. While many \gls{neurodivergent} 
developers possess high aptitude for logic and pattern recognition, they frequently encounter barriers related to \gls{executive_dysfunction}. 
This is an umbrella term for difficulty in mental processes that enable us to plan, focus attention, remember instructions, and juggle multiple tasks 
successfully. For a developer, \gls{executive_dysfunction} is not a lack of technical ability, but rather a barrier in three key areas: Inhibitory Control
 (filtering \gls{IDE} distractions), Working Memory (tracking complex syntax), and Cognitive Flexibility (switching between high-level logic and low-level code).

To provide a clearer lens for the objectives of this project, the below table outlines the specific strengths and environmental challenges associated with the conditions most prevalent in the developer community~\cite{sperry2005perceptions}~\cite{hotte2024strengths}~\cite{roitsch2019overview}.
\begin{table}[H]
    \centering
    \caption{\Gls{neurodivergent} Strengths and Environmental Challenges}
    \renewcommand{\arraystretch}{1.5} % Adds breathing room between rows
    \begin{tabularx}{\textwidth}{l >{\raggedright\arraybackslash}X >{\raggedright\arraybackslash}X}
        \toprule
        \textbf{Condition} & \textbf{Potential Strengths} & \textbf{Common Environmental Challenges} \\
        \midrule
        \gls{ASD} & Deep focus, systematic thinking, high integrity, exceptional memory for details, strong pattern recognition. & Sensory overload in noisy or bright environments, misinterpretation of direct communication style, difficulty with unspoken social rules. \\
        \addlinespace
        \gls{ADHD} & Creativity, hyperfocus on passions, adaptability, high energy, ability to multitask or rapidly switch contexts. & Distractibility in chaotic settings, struggles with time management and organisation leading to \gls{time-blindness} and issues with task initiation, impulsivity in structured environments.\\
        \addlinespace
        \gls{Dyslexia} & Strong visual-spatial reasoning, excellent big-picture thinking, creative problem-solving, strong verbal skills. & Difficulties with text-heavy tasks, timed reading assignments, and conventional written communication. \\
        \bottomrule
    \end{tabularx}
\end{table}

\subsection{Aims and Objectives}

\subsubsection{Project Aim}

The primary objective of this project is to design, implement, and evaluate an assistive VS Code plug-in that reduces extraneous cognitive load and mitigates the effects of
 \gls{executive_dysfunction} for \gls{neurodivergent} software developers. By introducing a cognitive prosthetic layer, the project seeks
  to determine if targeted \gls{UI} simplification and \gls{AI}-driven task scaffolding can measurably improve the developer experience.

\subsubsection{Research Objectives}

\textbf{Objective 1:} Interface Optimisation (The Minimalist Toggle)

Target: Develop a single ``Minimalist Toggle'' feature that simultaneously hides non-essential \gls{UI} elements, including the Activity Bar, Status Bar, Mini-map, and Breadcrumbs.


Justification: This directly addresses the issue of visual overload by reducing the number of stimuli competing for the developer's attention, thus lowering the Extraneous
 Cognitive Load and creating a more focused coding environment \cite{sweller1988cognitive}, reducing visual noise that leads to distraction and sensory overwhelm for autistic and \gls{ADHD} developers.

\textbf{Objective 2: }\gls{AI}-Driven Task Decomposition (The Task Architect)

Target: To implement an \gls{LLM}-integrated module that translates ambiguous programming goals into structured, editable checklists using a Chain of Thought (CoT) prompting strategy. This should be
implemented via the Open\gls{AI} \gls{API} that breaks down high-level requirements~\cite{ringel2019ai}.


Justification: This supports task initiation and reduces the energy required to overcome task paralysis, a primary barrier for developers with \gls{ADHD} and \gls{ASD}~\cite{gama2025socio}.

\textbf{Objective 3:} Simplification of Technical Feedback (The Code Explainer)

Target: To build a diagnostic interceptor that re-interprets complex compiler errors into plain-language, bulleted instructions~\cite{rocha2025conceptual}. Limit output to a maximum of three bullet points using direct instructions style.


Justification: This minimizes the cognitive demand of interpreting dense jargon, preventing rabbit holes and supporting the literal communication needs identified in the Double Empathy
 problem. By simplifying the feedback loop, developers can maintain focus on the core problem rather than getting lost in peripheral details. This is particularly beneficial for those with \gls{ASD},
  who may struggle with implied logic and prefer clear, direct instructions. Additionally, it can help mitigate the hyperfocus trap by providing concise guidance that keeps the developer on track 
  without overwhelming them with information.

\textbf{Objective 4: }Visual Accessibility Customisation (The Visual Styler)

Target: To create at least two visual presets, such as ``Soft High Contrast'', that adhere to \gls{WCAG} 2.1, set of international standards designed to make web content more accessible to people with disabilities it can also help with visual processing, eye strain and enhance readability~\cite{bengtsson2025web}.

Justification: Addresses the sensory sensitivities and colour preferences reported by \gls{neurodivergent} developers in research surveys~\cite{rocha2025conceptual}.

\textbf{Objective 5:} Empirical Evaluation of Cognitive Impact

Target: To conduct a comparative study using the \gls{NASATLX} and the \gls{SUS} to validate the plugin's efficacy against a standard \gls{IDE}.

Justification: This provides a standardized scientific metric to confirm if the tool successfully reduces mental demand and frustration for the target demographic.

\subsection{Dissertation Structure}

This dissertation is organized into five sections, following a logical progression from theoretical grounding to empirical validation:
\begin{itemize}
  \item \textbf{Section 2: Literature Review} I will evaluate the foundational psychological frameworks of Cognitive Load Theory (CLT) and Multimedia Learning (CTML). This provides a theoretical foundation for other research into the current state of \gls{neurodivergent} 
  software engineering as well as the shift toward a more social model of disability rather than a medical model.
  \item \textbf{Section 3: Design and Implementation} Details the user requirements and the three-layer system architecture. This chapter explains the technical
   execution of ``Minimalist Toggle'' and ``Task Architect'' as well as the visual presets and styling tools. I will also discuss issues encountered during development and how they were resolved, such as the design of the UI and issues with the AI integration and prompt generation.
  \item \textbf{Section 4: Results and Discussion} Presents the quantitative and qualitative findings from the user study ($N=5$). It analyzes the 82\% reduction in perceived workload measured by the NASA-TLX 
  and evaluates the tool`s excellent usability rating via the System Usability Scale (SUS). This section also includes a critical discussion of the results, exploring the implications for future assistive 
  technology in software development and acknowledging the limitations of the study, such as the small sample size and potential biases in participant selection.
  \item \textbf{Section 5: Conclusion} Summarises the core findings, acknowledges research limitations, and proposes specific avenues for future work with larger, more diverse participant pools.
\end{itemize}

\newpage